\section{Code Translation Pipeline}

The code translation pipeline used in this thesis consists of several sequential
steps, beginning with data preparation and ending with evaluation of the
generated code. The same overall pipeline is applied to both the base and
fine-tuned models in order to ensure fair comparison.

The first step is data loading and preprocessing. For the fine-tuning stage, the
XLCoST dataset is loaded and filtered to retain only Java-to-C++ translation
pairs. The Java code is used as the input sequence, while the corresponding C++
code serves as the target sequence. The selected code pairs are then converted
into instruction-based prompts suitable for causal language models.

The prompt format follows a simple instruction style. Each example contains a
task description such as ``Translate the following Java code into C++'' followed
by the original Java source code. During training, the expected C++ translation
is appended to the prompt so that the model can learn the mapping between the
input and output languages. Similar prompt-based approaches are commonly used in
modern code translation systems because they align well with instruction-tuned
large language models and simplify the translation task
\cite{aljagthami2025prompt, qwen_coder}.

After preprocessing, the models generate C++ code from the Java input samples.
For the base models, this step is performed directly using the original
pre-trained checkpoint. For the fine-tuned models, the LoRA adapter weights are
loaded on top of the base model before inference. In both cases, the generated
output is post-processed to remove markdown formatting, explanations, or other
unwanted text outside the code block.

%szb: még nem olvastam, de már látom, hogy a 3.5 az evaluaton-ről ír. Össze kellene rántani a methodot: ugyanazokat a dolgokat ismételgeted különböző sectionökben. Pl. az evaluation-t már említetted korábban, itt is leírod, meg a 3.5 is felteszem arról szól. Elég 1x az elején említeni, hogy mit használsz, majd utána amikor eljutsz oda a pipeline-ban, akkor leírni a konkrétumokat
The generated C++ code is then evaluated in two different ways. First, the
translation is compared against the reference output using BLEU score. BLEU
measures token-level similarity between the generated code and the reference
translation, providing a fast and widely used indicator of translation quality.

Second, execution-based evaluation is performed using the CodeEditorBench
framework. CodeEditorBench is designed specifically for realistic code editing
tasks, including translation, debugging, and refactoring. %szb: fentebb már bemutattad a benchmarkot, nem kell újra
Instead of measuring only textual similarity, the framework compiles and
executes the generated code against predefined test cases to determine whether
the translated code behaves correctly.

The CodeEditorBench evaluation was executed inside a Docker container using the
official benchmark image. Docker-based evaluation ensures that all models are
tested under the same environment, including compiler versions, runtime
dependencies, and execution settings. %szb: ne már, már harmadjára olvasom ugyanazt a docker-ről......
This is important because even small differences in environment configuration
can affect whether generated code compiles or executes correctly.

The full translation pipeline can therefore be summarized as follows: first, the %szb: nem szükséges a summary, effektíve 20adjára is megismétled ugyanazt. Legyen lényegretörőbb a szöveg úgy hogy az olvasó ne vesszen el az ismétlések közepette a szövegben, és akkor nem kell summary
Java source code is loaded and formatted into an instruction-based prompt.
Second, the model generates a C++ translation. Third, the generated code is
cleaned and post-processed. Finally, the resulting translation is evaluated
using both BLEU score and execution-based testing through CodeEditorBench.

This pipeline enables direct comparison between different models and between
base and fine-tuned variants. It also makes it possible to analyze whether
improvements in token-level similarity correspond to actual improvements in
functional correctness. %szb: felesleges